
When practitioners ask which training examples most influenced a given prediction, different attribution methods give substantially different answers. In experiments conducted for this work, four widely-used data attribution methods---TRAK, TracIn, Influence Functions (IF), and FastIF---were compared under controlled computational budgets. The methods identified largely non-overlapping sets of influential training examples, with pairwise Jaccard similarity of top-50 rankings below 1\%.

Data attribution has applications in model debugging by identifying training examples causing failures, data valuation for marketplace applications, curriculum learning optimization, and machine learning auditing through influence tracing. Despite methodological advances from foundational influence functions to efficient approximations like TRAK and DataInf, practitioners lack systematic guidance for selecting among these methods.

\subsubsection{1.1 The Multi-Objective Nature of Attribution Quality}

The conventional approach treats data attribution as a single well-posed problem: given a model and a test example, identify the most influential training points. Methods are typically evaluated on a single metric, usually rank correlation with leave-one-out (LOO) retraining, and methods achieving higher correlation are considered better. This framing implicitly assumes that improving one quality dimension improves all others.

This work examines this assumption. The analysis reveals that attribution quality encompasses at least two distinct aspects measured in these experiments:

\begin{itemize}
\item \textbf{Rank preservation (rho\_r):} How well does the method preserve the ordering of influential examples?
\item \textbf{Magnitude fidelity (rho\_m):} How accurately does the method estimate the magnitude of influence?
\end{itemize}

These dimensions serve different downstream applications. Rank preservation matters for identifying top-k influential examples, while magnitude fidelity matters for data pricing or weighting decisions.

\subsubsection{1.2 From Single-Metric to Multi-Objective Evaluation}

The experiments observe that methods face trade-offs between these quality dimensions. Specifically, Influence Functions achieve higher rank preservation while FastIF achieves higher magnitude fidelity at matched computational budgets. This pattern persists across all five tested compute budgets (10, 25, 50, 75, 100 gradient-equivalent operations).

What causes this decoupling? The hypothesis tested is that the answer lies in the geometry of optimization. In convex settings like logistic regression, the loss landscape has a unique global minimum, and the Hessian structure may ensure that all quality metrics move together. In non-convex deep networks, multiple local minima exist, and different approximation methods (random projections for TRAK, HVP iterations for IF, gradient similarity for TracIn) may navigate this landscape differently.

\subsubsection{1.3 Contributions}

This work makes the following contributions:

\begin{enumerate}
\item \textbf{Pareto characterization of data attribution methods.} The experiments demonstrate that finite-compute attribution exhibits trade-offs across quality dimensions, with IF vs FastIF showing metric crossings at all five tested compute budgets.
\end{enumerate}

\begin{enumerate}
\item \textbf{Comparison of convex and non-convex settings.} The experiments establish that metric coupling persists in convex settings (partial correlation at least 0.99) but breaks down in deep networks (R-squared drops to 0.034), consistent with the hypothesis that trade-offs arise from non-convex geometry rather than approximation quality alone.
\end{enumerate}

\begin{enumerate}
\item \textbf{Quantitative evidence of practical disagreement.} The experiments show that different attribution methods identify largely different influential training examples, with top-50 Jaccard similarity as low as 0.0024.
\end{enumerate}

\begin{enumerate}
\item \textbf{Unified evaluation framework.} The experiments introduce compute-normalized comparison via gradient-equivalent operations, enabling cross-method evaluation at matched computational cost.
\end{enumerate}
